\subsection{Invariant subspaces under a compact operator}

The following theorem is to be compared with Schur's lemma (A, VIII, p. 43, cor. and V, p. 386, prop. 6) and with cor. 4 of I, p. 26.

THEOREM 3. --- Let E be a separated locally convex space over C and let A be a subset of \(\mathcal{L}(\mathrm{E})\) . Let us make the following hypotheses:

\begin{enumerate}
\def\labelenumi{(\roman{enumi})}
\item
  There exists no closed vector subspace of E, distinct from \(\{0\}\) and E, which is stable under A;
\item
  The set A contains a nonzero compact endomorphism. Then the commutant of A is reduced to scalar multiples of the identity.
\end{enumerate}

By replacing A with the unital subalgebra of \(\mathcal{L}(\mathrm{E})\) generated by A, we reduce to the case where A is a unital subalgebra of \(\mathcal{L}(\mathrm{E})\) . Let then \(h\) be a nonzero compact endomorphism of E belonging to A.

Lemma 1. --- There exists an element a of A such that the kernel of \(ha - 1_{E}\) is not reduced to 0.

There exists an element \(x_{0}\) of E such that \(h(x_{0}) \neq 0\) . Let V be a closed neighborhood of \(h(x_{0})\) not containing 0. Since the endomorphism h is compact, one can choose an open convex neighborhood U of \(x_{0}\) such that \(h(\mathrm{U})\) is a relatively compact subset of V. The closure C of \(h(\mathrm{U})\) is therefore a compact convex subset of E; since it is contained in V, it does not contain 0.

Let \(y\) be a point of C. Denote by \(Ay\) the set of images of \(y\) by the elements of A. It is a vector subspace of E stable under A; it is nonzero since it contains \(y\) . By hypothesis (i) of the theorem, the set \(Ay\) is dense in E. There therefore exists an element \(b\) of A such that \(b(y) \in U\) . Since the set C is compact, there exists a finite family \((b_1, \ldots, b_n)\) of elements of A such that the sets \(b_i^{-1}(U)\) cover C. There exists a continuous partition of unity \((\varphi_1, \ldots, \varphi_n)\) on C subordinate to this covering (TG, IX, p. 47, th. 3). Let us define a map \(f: C \to E\) by setting

\[
f (x) = \sum_ {i = 1} ^ {n} \varphi_ {i} (x) b _ {i} (x)
\]

for all \(x \in \mathrm{C}\) ; it is a continuous map. Since U is convex and \(b_i(x)\) belongs to U when \(\varphi_i(x)\) is not zero, we have \(f(\mathrm{C}) \subset \mathrm{U}\) and \(h(f(\mathrm{C})) \subset \mathrm{C}\) . By the Leray--Schauder theorem (th. 2 of III, p. 11), there exists an element \(x \in \mathrm{C}\) such that \(h(f(x)) = x\) . Then set

\[
a = \sum_ {i = 1} ^ {n} \varphi_ {i} (x) b _ {i}.
\]

One has \(h(a(x)) = h(f(x)) = x\) and x is nonzero since 0 does not belong to C, hence the kernel of \(ha - 1_{E}\) is nonzero.

Let us finish the proof of Theorem 3. Let \(a\) be an element of A such that the kernel T of \(ha - 1_{\mathrm{E}}\) is nonzero (lemma 1). Let \(i\) denote the canonical injection of T into E. The linear map \(i\) is equal to \(h \circ a \circ i\) , hence it is compact. Thus the dimension of T is finite (remark 3 of III, p. 2). Let \(u\) be an element of \(\mathcal{L}(\mathrm{E})\) which commutes with A. Since \(u\) commutes with \(h \circ a\) , we have \(u(\mathrm{T}) \subset \mathrm{T}\) . Since T is finite-dimensional and nonzero over the algebraically closed field \(\mathbf{C}\) , the endomorphism of T induced by \(u\) admits an eigenvalue \(\lambda\) .

Set \(\mathrm{F} = \mathrm{Ker}(u - \lambda1_{\mathrm{E}})\) . This is a closed nonzero vector subspace of E; it is stable under A, since u commutes with the elements of A. By hypothesis (i), we have F = E, whence \(u = \lambda1_{E}\) .

COROLLARY 1. --- Keeping the hypotheses of Theorem 3 and supposing moreover that the elements of A are pairwise commuting. Then E is one-dimensional over C.

By hypothesis (ii) of Theorem 3, the space E is nonzero. It therefore contains a vector subspace F of dimension 1. This subspace is closed since E is assumed separated. By th. 3, every element of A is a homothety, hence stabilizes F. It then follows from hypothesis (i) of th. 3 that F is equal to E.

COROLLARY 2 (Lomonosov). --- Let E be a separated locally convex space of dimension at least 2 over the field C and let u be an endomorphism of E. Suppose that there exists a compact endomorphism \(h \neq 0\) of E commuting with u. There then exists a closed vector subspace F of E, distinct from \(\{0\}\) and from E, such that \(u(F) \subset F\) .

Indeed, Corollary 1 shows that the set A = \{u, h\} cannot satisfy hypothesis (i) of Theorem 3.

COROLLARY 3. --- Let E be a separated locally convex space over the field C and let u be a compact endomorphism of E. If E has dimension at least 2, there exists a closed vector subspace F of E, distinct from \(\{0\}\) and from E, such that \(u(F) \subset F\) .

The case u = 0 is immediate. The case \(u \neq 0\) follows from cor. 2 by taking h = u.

